% Lösungen Einheit 2 von 3 – Zielfunktion aus Haupt- und Nebenbedingung (zusammenbau v0.1)
\einheitenkopf{Einheit 2 von 3 · Zielfunktion aus Haupt- und Nebenbedingung}

\erg{11}{a) $2a + b = 60$ \quad b) Hauptbedingung $A = a \cdot b$, Nebenbedingung $2a + 2b = 40$, also $b = 20 - a$ \quad c) $b = 30 - 2a$: für $a = 6$ ist $b = 18$, $A = 108$ m²; für $a = 10$ ist $b = 10$, $A = 100$ m² \quad d) $A = a \cdot b$ und $2a + b = 70$, also $b = 70 - 2a$ und $A(a) = a \cdot (70 - 2a) = 70a - 2a^2$ mit $0 < a < 35$ \quad e) Breite $2x$, Höhe $2f(x)$: $A(x) = 2x \cdot 2 \cdot (5 - 0{,}2x^2) = 20x - 0{,}8x^3$ \quad f) $A(u) = u \cdot f(u) = (u^2 + 3u) \cdot e^{-0{,}5u}$ für $u > 0$ \quad g) Rechter Winkel bei $(a | 0)$: Grundseite $a$, Höhe $f_k(a) = 2a \cdot e^{-ka}$, also $A = \tfrac{1}{2} \cdot a \cdot 2a \cdot e^{-ka} = a^2 \cdot e^{-ka}$.}
\erg{12}{a) Gesucht ist eine Funktion von $a$, also nach $b$ umstellen: $b = 50 - 2a$ und $A(a) = a \cdot (50 - 2a) = 50a - 2a^2$. \quad b) Nur eine Funktion einer Variablen lässt sich ableiten und über die Nullstellen der Ableitung untersuchen; die zweite Variable hängt über die Nebenbedingung von der ersten ab und ist nicht frei wählbar. \quad c) $\pi r^2 h = 1\,000$, also $h = \frac{1\,000}{\pi r^2}$; $O(r) = 2\pi r^2 + 2\pi r h = 2\pi r^2 + \frac{2\,000}{r}$ für $r > 0$; $O(5) = 50\pi + 400 \approx 557{,}1$ cm²}
